% notes/artefacts/execution-market-making/market-participants-trading-motives.tex
% Legacy file kept verbatim when it was migrated (notes/source-review-records.md, market design and liquidity / source review / optimal execution / execution control reviews): the whole of
% compendium/legacy/execution-market-making/01_foundations_of_trading_and_market_organization/02_market_participants_and_trading_motives.tex
% below its migrate-to line (migrate-to: execution-market-making, market-structure-execution-quality), with only the subsection heading unnumbered. The compendium entries
% state each claim with its locator; the original wording lives here. Treat everything here as unchecked own working.
% Destinations: Trading Game View, Long and Short Positions, Motive Segmentation, Multiple Trader Identities, Patience Defines the Liquidity Side, Buy-Side Motives, Instrument Choice Follows the Problem, Hedge-Vehicle Fit, Profit-Motivated Trader Segmentation, Speculation Requires Comparative Advantage, Trading Edge Segmentation, Noise and Liquidity Traders, Futile Trader Segmentation -> def:trader-motives; Skill and Luck Look Similar Ex Post, Past Profit Alone Does Not Identify Skill (proposition, conditions restated), the two corollaries with proofs -> res:zero-sum-skill-luck; Liquidity Supplier Segmentation, Dealers and Market Makers, Sell-Side Intermediaries, Broker Service Segmentation -> def:market-intermediaries; Information-Motivated Traders, Informed Trader Liquidity Role Segmentation -> def:informed-trader-types; Arbitrageur Economic Function -> def:arbitrage; Block-Trading Role Segmentation, Block Dealers vs Block Brokers, Block Trader Skill Segmentation -> def:block-trade (routing exception); Broker Agency Conflict -> def:broker-agency-conflicts; Institutional Investors and HFTs, HFT Identification Markers, HFT as Fragmentation Link, Blurring of Roles -> def:hft-strategies (routing exception); Tighter Spread Different Cost, Investor Size Asymmetry -> res:hft-spread-investor-cost (routing exception). From Harris chaps. 3, 7, 8, 15, 19 and 22, Lehalle and Laruelle's introduction.
% Compile inside notes/artefacts/execution-market-making/ with  latexmk -pdf market-participants-trading-motives.tex
\PassOptionsToPackage{table}{xcolor}
\documentclass[a4paper,oneside,openany]{book}
% Your own packages, macros and environments. The learning guide's entry environments,
% \citep and \qfsource are in house.tex, which compendium.tex loads after this file.

% -----------------------------
% Document Class & Encoding
% -----------------------------
\usepackage[utf8]{inputenc}  % Input encoding
\usepackage[T1]{fontenc}     % Font encoding

% -----------------------------
% Page Layout & Geometry
% -----------------------------
\usepackage[margin=2cm,foot=1cm]{geometry}

% -----------------------------
% Header & Footer Management
% -----------------------------
\usepackage{fancyhdr}
\pagestyle{fancy}

% -----------------------------
% Text Formatting & Spacing
% -----------------------------
\usepackage{ragged2e}
\usepackage[nodisplayskipstretch]{setspace}
\setstretch{1}  % Line spacing

% -----------------------------
% Listings for Code
% -----------------------------
\usepackage{listings}
\usepackage{xcolor}

% -----------------------------
% Tables, Figures & Graphics
% -----------------------------
\usepackage{tabularx}
\usepackage{multirow}
\usepackage{pdflscape}
\usepackage{graphicx}

% -----------------------------
% Algorithms
% -----------------------------
\usepackage{algorithm}       % For floating algorithm environment
\usepackage{algpseudocode}   % For pseudo-code inside algorithms

% -----------------------------
% Math Packages
% -----------------------------
\usepackage{amsmath,amssymb,amsthm,mathtools,bm}
\usepackage{mathrsfs}        % Script fonts
\usepackage{tikz-cd}         % Commutative diagrams
\usepackage{rotating}        % Rotated objects
\usepackage{float}           % Float placement
\usetikzlibrary{calc, arrows, positioning, arrows.meta}

% -----------------------------
% TikZ and its Libraries
% -----------------------------
\usepackage{tikz}
\usetikzlibrary{positioning} % Added to support "below=of" syntax
\usepackage{forest}          % For advanced tree diagrams
\usepackage{tikz-qtree}      % For vertical tree diagrams

% -----------------------------
% Miscellaneous Packages
% -----------------------------
\usepackage{xparse}          % For defining robust commands
\usepackage{enumitem}        % For customizing list layouts
\usepackage{makecell}        % For enhanced table cells
\usepackage{comment}		 % For commenting out sections

%%%%%%%%%%%%%%%%%%%%%%%%%%%%%%%%%%%%%%%%%%%%%%%%%%%%%%%%%%%%%%%%%%%%%%%%%%%%%%%
% Custom Math Commands and Shortcuts
%%%%%%%%%%%%%%%%%%%%%%%%%%%%%%%%%%%%%%%%%%%%%%%%%%%%%%%%%%%%%%%%%%%%%%%%%%%%%%%

% Matrix and Vector Formatting
\newcommand{\matr}[1]{\mathbf{#1}}      % For undergraduate algebra (alternative: \bm{#1})
\newcommand{\vect}[1]{\bm{#1}}

% Standard Number Sets
\newcommand{\N}{\mathbb{N}}
\newcommand{\R}{\mathbb{R}}
\newcommand{\Q}{\mathbb{Q}}
\newcommand{\C}{\mathbb{C}}
\newcommand{\Z}{\mathbb{Z}}

% Bold variables (customize as needed)
\newcommand{\x}{\mathbf{x}}
\newcommand{\F}{\mathbf{F}}
\newcommand{\f}{\mathbf{f}}
\newcommand{\y}{\mathbf{y}}
\renewcommand{\a}{\mathbf{a}}
\renewcommand{\b}{\mathbf{b}}
\renewcommand{\c}{\mathbf{c}}
\newcommand{\h}{\mathbf{h}}
\newcommand{\g}{\mathbf{g}}
\newcommand{\z}{\mathbf{z}}
\newcommand{\ze}{\mathbf{0}}

% Norms, Absolute Values, and Bracketing
\newcommand{\norm}[1]{\left\lVert#1\right\rVert}
\newcommand{\abs}[1]{\left\lvert#1\right\rvert}
\newcommand{\brk}[1]{\left[#1\right]}
\newcommand{\brc}[1]{\left\{#1\right\}}
\newcommand{\paren}[1]{\left(#1\right)}
\newcommand{\normop}[1]{\left\lVert#1\right\rVert_\text{op}}

% Miscellaneous Math Shortcuts
\newcommand{\LL}{\mathcal{L}}
\newcommand{\uni}{\overset{\text{uni}}{\to}}
\DeclareMathOperator{\diam}{diam}
\newcommand{\Prr}[1]{\text{Pr}\left(#1\right)}
\newcommand{\ceil}[1]{\lceil #1 \rceil}
\newcommand{\floor}[1]{\lfloor #1 \rfloor}
\newcommand\disteq{\stackrel{\textnormal{dist}}{=}}

% Evaluation command with options
\NewDocumentCommand{\evalat}{sO{\big}mm}{%
  \IfBooleanTF{#1}
    {\mleft. #3 \mright|_{#4}}
    {#3#2|_{#4}}%
}

%%%%%%%%%%%%%%%%%%%%%%%%%%%%%%%%%%%%%%%%%%%%%%%%%%%%%%%%%%%%%%%%%%%%%%%%%%%%%%%
% Color Definitions
%%%%%%%%%%%%%%%%%%%%%%%%%%%%%%%%%%%%%%%%%%%%%%%%%%%%%%%%%%%%%%%%%%%%%%%%%%%%%%%

\definecolor{dartmouthgreen}{rgb}{0.05, 0.5, 0.06}
\definecolor{egyptianblue}{rgb}{0.06, 0.2, 0.65}
\definecolor{dukeblue}{rgb}{0.0, 0.0, 0.61}
\definecolor{jazzberryjam}{rgb}{0.65, 0.04, 0.37}
\definecolor{magenta}{HTML}{EC008C}
\definecolor{darkmagenta}{rgb}{0.55, 0.0, 0.55}
\definecolor{deeppink}{rgb}{1.0, 0.08, 0.58}
\definecolor{codegreen}{rgb}{0,0.6,0}
\definecolor{codegray}{rgb}{0.5,0.5,0.5}
\definecolor{codepurple}{rgb}{0.58,0,0.82}
\definecolor{backcolour}{rgb}{0.95,0.95,0.92}

%%%%%%%%%%%%%%%%%%%%%%%%%%%%%%%%%%%%%%%%%%%%%%%%%%%%%%%%%%%%%%%%%%%%%%%%%%%%%%%
% Custom Shortcuts and Emphasis Commands
%%%%%%%%%%%%%%%%%%%%%%%%%%%%%%%%%%%%%%%%%%%%%%%%%%%%%%%%%%%%%%%%%%%%%%%%%%%%%%%

\newcommand{\isomorp}{\xrightarrow{\sim}}
\newcommand{\hlt}[1]{\textit{{\color{blue}#1}}}
\newcommand{\simpt}[1]{\textit{{\color{deeppink}#1}}}
\newcommand{\impt}[1]{\textit{{\color{red}#1}}}
\newcommand{\gcds}[1]{\textnormal{gcd}#1}
\newcommand{\mins}[1]{\textnormal{min}#1}
\newcommand{\maxs}[1]{\textnormal{max}#1}
\newcommand{\lcms}[1]{\textnormal{lcm}#1}
\newcommand{\degs}[1]{\textnormal{deg}#1}
\newcommand{\exps}[1]{\textnormal{exp}#1}
\newcommand{\tors}[1]{\textnormal{Tor}#1}
\newcommand{\homs}[1]{\textnormal{Hom}#1}
\newcommand{\anns}[1]{\textnormal{Ann}#1}
\newcommand{\defeq}{\vcentcolon=}
\DeclareMathOperator{\sign}{sign}

%%%%%%%%%%%%%%%%%%%%%%%%%%%%%%%%%%%%%%%%%%%%%%%%%%%%%%%%%%%%%%%%%%%%%%%%%%%%%%%
% Theorem and Definition Environments
%%%%%%%%%%%%%%%%%%%%%%%%%%%%%%%%%%%%%%%%%%%%%%%%%%%%%%%%%%%%%%%%%%%%%%%%%%%%%%%
\theoremstyle{plain}
\newtheorem{theorem}{Theorem}[section]
\newtheorem{corollary}[theorem]{Corollary}
\newtheorem{lemma}[theorem]{Lemma}
\newtheorem{proposition}[theorem]{Proposition}
\newtheorem{qnbank}[theorem]{Question Bank}
\let\oldqnbank\qnbank
\renewcommand{\qnbank}{\oldqnbank\normalfont}
\newtheorem{process}[theorem]{Process}
\let\oldprocess\process
\renewcommand{\process}{\oldprocess\normalfont}
\newtheorem{method}[theorem]{Method}
\let\oldmethod\method
\renewcommand{\method}{\oldmethod\normalfont}
\newtheorem{definition}[theorem]{Definition}
\let\olddefinition\definition
\renewcommand{\definition}{\olddefinition\normalfont}
\newtheorem{example}[theorem]{Example}
\let\oldexample\example
\renewcommand{\example}{\oldexample\normalfont}
\newtheorem{remark}[theorem]{Remark}
\let\oldremark\remark
\renewcommand{\remark}{\oldremark\normalfont}
\newtheorem{notation}[theorem]{Notation}
\let\oldnotation\notation
\renewcommand{\notation}{\oldnotation\normalfont}
\newtheorem{principle}[theorem]{Principle}
\let\oldprinciple\principle
\renewcommand{\principle}{\oldprinciple\normalfont}
\newtheorem{mechanism}[theorem]{Mechanism}
\let\oldmechanism\mechanism
\renewcommand{\mechanism}{\oldmechanism\normalfont}
\newtheorem{taxonomy}[theorem]{Taxonomy}
\let\oldtaxonomy\taxonomy
\renewcommand{\taxonomy}{\oldtaxonomy\normalfont}
\newtheorem{model}[theorem]{Model}
\let\oldmodel\model
\renewcommand{\model}{\oldmodel\normalfont}
\newtheorem{assumption}[theorem]{Assumption}
\let\oldassumption\assumption
\renewcommand{\assumption}{\oldassumption\normalfont}
\newtheorem{casestudy}[theorem]{Case Study}
\let\oldcasestudy\casestudy
\renewcommand{\casestudy}{\oldcasestudy\normalfont}
\newtheorem{empiricalfinding}[theorem]{Empirical Finding}
\let\oldempiricalfinding\empiricalfinding
\renewcommand{\empiricalfinding}{\oldempiricalfinding\normalfont}
\newtheorem{practitionernote}[theorem]{Practitioner Note}
\let\oldpractitionernote\practitionernote
\renewcommand{\practitionernote}{\oldpractitionernote\normalfont}
\newtheorem{caveat}[theorem]{Caveat}
\let\oldcaveat\caveat
\renewcommand{\caveat}{\oldcaveat\normalfont}
\newtheorem{marketrule}[theorem]{Rule}
\let\oldmarketrule\marketrule
\renewcommand{\marketrule}{\oldmarketrule\normalfont}
\newtheorem{takeaway}[theorem]{Takeaway}
\let\oldtakeaway\takeaway
\renewcommand{\takeaway}{\oldtakeaway\normalfont}

%%%%%%%%%%%%%%%%%%%%%%%%%%%%%%%%%%%%%%%%%%%%%%%%%%%%%%%%%%%%%%%%%%%%%%%%%%%%%%%
% Breakable Algorithm Environment (Algorithms spanning multiple pages)
%%%%%%%%%%%%%%%%%%%%%%%%%%%%%%%%%%%%%%%%%%%%%%%%%%%%%%%%%%%%%%%%%%%%%%%%%%%%%%%

\makeatletter
\newenvironment{breakablealgorithm}
 {%
  \begin{center}
    \refstepcounter{algorithm} % Increment algorithm counter
    \hrule height.8pt depth0pt \kern2pt
    \renewcommand{\caption}[2][\relax]{%
      {\raggedright\textbf{Algorithm \thealgorithm:} ##2\par}%
      \ifx\relax##1\relax
        \addcontentsline{loa}{algorithm}{\protect\numberline{\thealgorithm}##2}%
      \else
        \addcontentsline{loa}{algorithm}{\protect\numberline{\thealgorithm}##1}%
      \fi
      \kern2pt\hrule\kern2pt
    }%
 }{%
    \kern2pt\hrule
  \end{center}
 }
\makeatother

%%%%%%%%%%%%%%%%%%%%%%%%%%%%%%%%%%%%%%%%%%%%%%%%%%%%%%%%%%%%%%%%%%%%%%%%%%%%%%%
% Code Listings Style
%%%%%%%%%%%%%%%%%%%%%%%%%%%%%%%%%%%%%%%%%%%%%%%%%%%%%%%%%%%%%%%%%%%%%%%%%%%%%%%

\lstdefinestyle{codestyle}{
  commentstyle=\color{codegreen},
  keywordstyle=\color{magenta},
  numberstyle=\tiny\color{codegray},
  stringstyle=\color{codepurple},
  basicstyle=\ttfamily\normalsize,
  breakatwhitespace=false,
  breaklines=true,
  captionpos=b,
  keepspaces=true,
  numbers=left,
  numbersep=10pt,
  frame=single,
  showspaces=false,
  showstringspaces=false,
  showtabs=false,
  tabsize=2
}
\lstset{style=codestyle}

% -----------------------------
% Hyperlinks (loaded last: loaded before float and algorithm it duplicated figure anchors)
% -----------------------------
\usepackage{hyperref}
\hypersetup{
  colorlinks,
  linkcolor={blue!80!black},
  citecolor={green!80!black},
  urlcolor={red!80!black}
}

\usepackage[T1]{fontenc}
\usepackage{amsmath,amssymb,booktabs,graphicx,hyperref}
\hypersetup{colorlinks}
\setlength{\emergencystretch}{3em}   % lets a line holding a long key break
\newcounter{qfentry}[chapter]
\renewcommand{\theqfentry}{\thechapter.\arabic{qfentry}}
% \qftype{type}{Name} declares the environment qf<type>; its title wraps like a paragraph.
\newcommand{\qftype}[2]{\newenvironment{qf#1}[1][]{\par\addvspace{\medskipamount}%
  \refstepcounter{qfentry}{\raggedright\noindent\bfseries #2 \theqfentry
  \if\relax\detokenize{##1}\relax\else\quad ##1\fi\par}\nopagebreak}%
  {\par\addvspace{\medskipamount}}}
\qftype{definition}{Definition}\qftype{convention}{Convention}\qftype{result}{Result}
\qftype{model}{Model}\qftype{method}{Method}\qftype{mechanism}{Mechanism}
\qftype{evidence}{Evidence}\qftype{data}{Data}\qftype{rule}{Rule}
\qftype{pitfall}{Pitfall}\qftype{example}{Example}
\newcommand{\qffield}[1]{\par\textbf{#1.}\enspace}   % starts one of an entry's fields
% \citep[locator]{key} links the key to its line in the sources if it has one,
% and otherwise to its entry in the pinned catalogue PDF.
\newcommand{\qfcatalogue}{../guide/institutional_quantitative_finance_reading_list.pdf}
\NewDocumentCommand{\citep}{o o m}{[\IfValueT{#2}{#1 }%
  \ifcsname r@src:#3\endcsname\hyperref[src:#3]{\nolinkurl{#3}}%
  \else\href{\qfcatalogue\##3}{\nolinkurl{#3}}\fi
  \IfValueTF{#2}{, #2}{\IfValueT{#1}{, #1}}]}
% \qfsource{key}{reference} records one source outside the catalogue, in front/sources.tex.
\newcommand{\qfsource}[2]{\par\noindent\phantomsection\label{src:#1}\nolinkurl{#1}\quad #2\par}
% \qfsections{key, ...} inputs <key>/section.tex for each listed key that has one, in order.
\ExplSyntaxOn
\NewDocumentCommand{\qfsections}{m}
  {\clist_map_inline:nn {#1} {\InputIfFileExists{##1/section}{}{}}}
\ExplSyntaxOff

\pagestyle{plain}
\begin{document}
\chapter*{Legacy notes: market participants and trading motives}

% From compendium/legacy/execution-market-making/01_foundations_of_trading_and_market_organization/02_market_participants_and_trading_motives.tex (whole file, deleted 2026-10-16 once empty).
\subsection*{Market Participants and Trading Motives}

\begin{model} \hlt{Trading Game View}\\
A trading mechanism is a game in which participants submit orders under venue rules. The players include customers, brokers, dealers, market makers, and traders acting for their own account. Some demand liquidity, some supply liquidity, some carry information, and some process orders.\\
Intuition: a trader's best order depends on what others may know, how fast they can react, and how the venue ranks competing messages.
\end{model}

\begin{notation} \hlt{Long and Short Positions}\\
A trader with a \hlt{long position} owns the asset and benefits when its price rises. A trader with a \hlt{short position} has sold an asset not owned, or has equivalent negative exposure, and benefits when its price falls. Closing a short position is called \hlt{covering}.
\end{notation}

\begin{taxonomy} \hlt{Motive Segmentation}
\begin{enumerate}[label=\roman*.]
\setlength{\itemsep}{0pt}
\item \hlt{utilitarian}: solve an outside problem;
\item \hlt{profit-motivated}: expect trading profits above costs;
\item \hlt{futile}: expect profit without enough information, skill, discipline, or agency control.
\end{enumerate}
Motive is inferred from context, horizon, and exposure, not order side alone.
\end{taxonomy}

\begin{caveat} \hlt{Multiple Trader Identities}\\
Traders often combine motives. A pension fund may invest, hedge FX, and time factors; an HFT may make markets in one book and arbitrage another.
\end{caveat}

\begin{principle} \hlt{Patience Defines the Liquidity Side}\\
Impatient traders demand liquidity by initiating trades. Patient traders supply liquidity by waiting, posting, quoting, or making themselves discoverable.\\
The same institution can demand liquidity for one order and supply it for another.
\end{principle}

\begin{taxonomy} \hlt{Buy-Side Motives}\\
Buy-side traders use markets to solve problems outside the market itself:
\begin{enumerate}[label=\roman*.]
\setlength{\itemsep}{0pt}
\item investors move wealth across time;
\item borrowers move wealth from the future to the present;
\item hedgers reduce business or portfolio risk;
\item asset exchangers acquire assets they value more than what they tender;
\item gamblers trade for entertainment or payoff convexity rather than investment efficiency.
\end{enumerate}
These traders buy liquidity services rather than sell them.
\end{taxonomy}

\begin{principle} \hlt{Instrument Choice Follows the Problem}\\
Investors and borrowers care about after-cost return and funding access. Asset exchangers need the right asset. Hedgers need correlation, maturity, and liquidity. Gamblers care about payoff shape.
\end{principle}

\begin{taxonomy} \hlt{Hedge-Vehicle Fit}
\begin{enumerate}[label=\roman*.]
\setlength{\itemsep}{0pt}
\item \hlt{forwards}: customise asset, date, size, and delivery, but create search, credit, and exit frictions;
\item \hlt{futures}: standardise terms and clearing, improving liquidity while introducing basis and delivery mismatch;
\item \hlt{options}: cap downside for a premium, preserving convex upside or nonlinear exposure;
\item \hlt{swaps}: exchange cash-flow risks over time, especially interest-rate, currency, credit, or inflation exposures.
\end{enumerate}
The best hedge removes the intended risk without adding larger liquidity, basis, or counterparty risks.
\end{taxonomy}

\begin{taxonomy} \hlt{Information-Motivated Traders}\\
Informed traders trade when prices lag information. Value traders use fundamentals; news traders use recent information; arbitrageurs use relative mispricing.\\
Their orders improve prices only as the market absorbs the information.
\end{taxonomy}

\begin{taxonomy} \hlt{Informed Trader Liquidity Role Segmentation}
\begin{enumerate}[label=\roman*.]
\setlength{\itemsep}{0pt}
\item \hlt{value traders}: supply liquidity when uninformed demand pushes price away from fundamental value;
\item \hlt{news traders}: demand liquidity when new information has not yet been incorporated into price;
\item \hlt{arbitrageurs}: supply or demand liquidity across related instruments to close relative mispricing.
\end{enumerate}
The same trader can switch roles. The useful distinction is whether the order corrects stale prices, absorbs non-value flow, or links related prices.
\end{taxonomy}

\begin{mechanism} \hlt{Arbitrageur Economic Function}\\
Arbitrageurs trade relative value. They connect buyers and sellers across instruments, venues, maturities, currencies, or risk forms. Their activity enforces the law of one price only after costs, funding, hedge error, and execution risk are small enough for the trade to be worth doing.
\end{mechanism}

\begin{taxonomy} \hlt{Liquidity Supplier Segmentation}
\begin{enumerate}[label=\roman*.]
\setlength{\itemsep}{0pt}
\item \hlt{market makers}: supply immediacy for small to moderate size at narrow spreads, but lose to informed flow and adverse inventory moves;
\item \hlt{block dealers}: supply depth to known clients by warehousing risk, but lose when client information or unwind cost is misread;
\item \hlt{value traders}: supply depth and resiliency when price deviates from fundamental value, but lose to better-informed value traders or financing pressure;
\item \hlt{precommitted traders}: post aggressive limit orders to lower the cost of trades they already want to do, but risk non-execution or adverse selection;
\item \hlt{arbitrageurs}: move liquidity across related markets, but lose to faster arbitrageurs, legging risk, funding shocks, and stale hedges.
\end{enumerate}
Treating all passive flow as ``market making'' hides the source of depth.
\end{taxonomy}

\begin{taxonomy} \hlt{Profit-Motivated Trader Segmentation}
\begin{enumerate}[label=\roman*.]
\setlength{\itemsep}{0pt}
\item \hlt{informed speculators}: profit when prices move toward information they already understand;
\item \hlt{order anticipators}: profit by predicting trades other participants must make;
\item \hlt{bluffers and manipulators}: profit only if others misread false signals;
\item \hlt{dealers}: sell immediacy and inventory capacity rather than directional views.
\end{enumerate}
Informed speculators aid price discovery; dealers supply liquidity. Middle categories mainly redistribute surplus.
\end{taxonomy}

\begin{principle} \hlt{Speculation Requires Comparative Advantage}\\
Profitable speculation requires more than being intelligent or informed in isolation. The trader must have an advantage over the marginal counterparty after spread, fees, financing, market impact, and operational error.
\end{principle}

\begin{taxonomy} \hlt{Trading Edge Segmentation}
\begin{enumerate}[label=\roman*.]
\setlength{\itemsep}{0pt}
\item \hlt{information}: faster, cleaner, or more exclusive interpretation of value-relevant facts;
\item \hlt{model}: better mapping from data to price, hedge, or risk;
\item \hlt{execution}: lower leakage, better routing, queue skill, or lower impact;
\item \hlt{risk bearing}: ability to hold inventory or convergence risk when others cannot;
\item \hlt{discipline}: sizing, stopping, and avoiding trades outside the edge.
\end{enumerate}
An edge that cannot survive trading costs is not a trading edge.
\end{taxonomy}

\begin{caveat} \hlt{Skill and Luck Look Similar Ex Post}\\
Large realised profits may come from skill, favourable exposure, or simple luck.\\
A useful skill claim must explain why the decision was right before the outcome was known, why the edge should persist, and why it is not already competed away.
\end{caveat}

\begin{proposition} \hlt{Past Profit Alone Does Not Identify Skill}\\
Let future excess profit be $\Pi_{t+1}=S_{t+1}+L_{t+1}$, where $S$ is persistent skill and $L$ is zero-mean luck.\\
If $S_{t+1}$ is independent of past realised profits conditional on current information, then past profit has no incremental predictive value for $\Pi_{t+1}$.
\begin{proof}
Conditional on current information $\mathcal{F}_t$, $\mathbb{E}[\Pi_{t+1}\mid \mathcal{F}_t,\Pi_t]=\mathbb{E}[S_{t+1}\mid \mathcal{F}_t]+\mathbb{E}[L_{t+1}\mid \mathcal{F}_t]=\mathbb{E}[\Pi_{t+1}\mid \mathcal{F}_t]$.
\end{proof}
\end{proposition}

\begin{corollary} \hlt{Profit-Motivated Flow Needs Non-Profit Motives}\\
Speculators and dealers cannot all profit by trading only with each other. Their gains require traders who pay for immediacy, risk transfer, information, or entertainment.
\begin{proof}
Let $G_i$ be trader $i$'s gross trading gain in a closed set of profit-motivated traders, and let $C_i\ge0$ be trading costs. By zero-sum accounting, $\sum_i G_i=0$, so aggregate net profit is $\sum_i (G_i-C_i)=-\sum_i C_i\le0$.\\
Thus the group cannot earn positive aggregate net trading profit from itself; positive profits require counterparties with offsetting trading losses or non-trading service value.
\end{proof}
\end{corollary}

\begin{corollary} \hlt{Turnover Drag for Active Management}\\
Let $R_M$ be the gross market return, $a_i$ active manager $i$'s gross active excess return, $c_i\ge0$ the manager's trading and fee drag, and $w_i$ value weights. If active managers aggregate to the market portfolio, $\sum_i w_i a_i=0$, so aggregate net active excess return is $\sum_i w_i(a_i-c_i)=-\sum_i w_i c_i\le0$.\\
A single active manager beats a passive fund with drag $c_P$ only if $a_i>c_i-c_P$.
\begin{proof}
Active deviations clear against each other in aggregate, giving $\sum_i w_i a_i=0$. Subtracting costs gives the aggregate inequality. Relative to passive exposure, manager $i$ must cover the incremental drag $c_i-c_P$.
\end{proof}
\end{corollary}

\begin{definition} \hlt{Noise and Liquidity Traders}\\
Noise traders trade for non-value reasons: rebalancing, cash needs, index changes, behavioural demand, or errors. They mask informed flow and give liquidity suppliers less toxic orders.
\end{definition}

\begin{taxonomy} \hlt{Futile Trader Segmentation}
\begin{enumerate}[label=\roman*.]
\setlength{\itemsep}{0pt}
\item \hlt{pseudo-informed}: trade stale or already-public information as if it were private;
\item \hlt{inefficient}: have a plausible strategy but execute, size, or control risk poorly;
\item \hlt{victimised}: rely on advisers, brokers, or agents whose incentives are misaligned;
\item \hlt{rogue}: hide losses or risks because internal controls fail.
\end{enumerate}
Futile traders matter because their losses can fund informed traders and dealers, while their distress can still damage clients, firms, or market confidence.
\end{taxonomy}

\begin{mechanism} \hlt{Dealers and Market Makers}\\
Dealers trade for their own account. Market makers quote prices or otherwise stand ready to trade, often holding inventory only until an offsetting order arrives. Their compensation comes from spreads, rebates, information, or client relationships; their risks are inventory, adverse selection, funding, and operational failure.
They are paid for immediacy: the client trades now, while the market maker manages the uncertainty of unwinding later.
\end{mechanism}

\begin{taxonomy} \hlt{Sell-Side Intermediaries}\\
Brokers arrange trades for clients without intending to own the position. Dealers trade as principal. Broker-dealers combine both functions, sometimes arranging agency trades and sometimes committing their own capital. The sell side exists because buy-side traders pay for search, immediacy, execution expertise, clearing access, and balance-sheet intermediation.
\end{taxonomy}

\begin{taxonomy} \hlt{Broker Service Segmentation}
\begin{enumerate}[label=\roman*.]
\setlength{\itemsep}{0pt}
\item \hlt{search}: find counterparties, dealers, venues, or natural block liquidity;
\item \hlt{representation}: submit, monitor, and negotiate client orders;
\item \hlt{access}: connect clients to exchanges, dealers, ECNs, and clearing members;
\item \hlt{post-trade}: manage confirmation, clearing, settlement, custody, margin, and records;
\item \hlt{information control}: expose enough order interest to trade without revealing the full parent order.
\end{enumerate}
\end{taxonomy}

\begin{taxonomy} \hlt{Block-Trading Role Segmentation}
\begin{enumerate}[label=\roman*.]
\setlength{\itemsep}{0pt}
\item \hlt{block initiator}: trader with a parent order too large for normal displayed liquidity;
\item \hlt{block liquidity supplier}: trader, dealer, or institution willing to fill the block;
\item \hlt{block dealer}: principal that facilitates the block and then unwinds inventory;
\item \hlt{block broker}: agent that assembles counterparties without intending to own the risk;
\item \hlt{upstairs market}: relationship and telephone market used to find size without public exposure.
\end{enumerate}
Block trading is mostly a search and information-control problem, not a different asset.
\end{taxonomy}

\begin{principle} \hlt{Block Dealers vs Block Brokers}\\
Dealers are useful when immediacy matters and balance sheet can absorb the position.\\
Brokers are useful when the block is too large or information-sensitive for one principal, so assembling natural counterparties is cheaper than warehousing risk.\\
Broker-dealers can be credible when they are willing to participate in the trade they recommend.
\end{principle}

\begin{taxonomy} \hlt{Block Trader Skill Segmentation}
\begin{enumerate}[label=\roman*.]
\setlength{\itemsep}{0pt}
\item \hlt{latent demand discovery}: know who might trade if contacted;
\item \hlt{motive audit}: infer whether the initiator is uninformed, informed, honest, or strategic;
\item \hlt{exposure control}: reveal size only to counterparties who will not trade ahead;
\item \hlt{transaction-cost analysis}: price facilitation, unwind risk, and likely market impact;
\item \hlt{reputation management}: make future counterparties willing to participate.
\end{enumerate}
\end{taxonomy}

\begin{caveat} \hlt{Broker Agency Conflict}\\
Brokers are agents, but commissions, soft dollars, payments for order flow, inventory, and affiliated dealers can pull routing decisions away from pure client execution quality.
\end{caveat}

\begin{practitionernote} \hlt{Institutional Investors and HFTs}\\
Institutional investors often trade large parent orders that must be hidden, sliced, or negotiated to reduce market impact. High-frequency traders use speed, automation, and short holding periods; they may act as market makers, arbitrageurs, statistical traders, or order anticipators. Their role depends on strategy.
\end{practitionernote}

\begin{taxonomy} \hlt{HFT Identification Markers}
\begin{enumerate}[label=\roman*.]
\setlength{\itemsep}{0pt}
\item automated order generation, routing, execution, and cancellation;
\item low-latency infrastructure, co-location, and direct market-data feeds;
\item short holding periods and low overnight inventory;
\item high message counts relative to executed trades;
\item market making, cross-venue arbitrage, or very short-horizon statistical trading.
\end{enumerate}
These markers describe technology and horizon, not whether the strategy is beneficial.
\end{taxonomy}

\begin{mechanism} \hlt{HFT as Fragmentation Link}\\
Fast market makers connect fragmented venues by quoting related prices, arbitraging stale quotes, and hedging inventory across books. Their compensation is the spread, rebates, and short-lived information advantages; their cost is technology, inventory risk, and adverse selection. They act like moving bridges between books: useful when they keep prices aligned, fragile when those bridges are cancelled under stress.
\end{mechanism}

\begin{caveat} \hlt{Tighter Spread, Different Cost}\\
An HFT can improve the displayed spread by stepping inside an older quote, but investor execution cost may not fall by the same amount. Smaller trades, more prints, and higher routing complexity can shift fixed and reporting costs toward slower investors.
\end{caveat}

\begin{caveat} \hlt{Investor Size Asymmetry}\\
The uncertainty created by meeting a fast counterparty is not borne equally. A small investor sees high variance in a small number of trades; a large institution diversifies this uncertainty across many child orders but still pays aggregate spread, impact, and signalling costs.
\end{caveat}

\begin{caveat} \hlt{Blurring of Roles}\\
Fragmentation blurs the distinction between liquidity takers and providers. The same agent may consume liquidity in one venue, supply it in another, and intermediate price formation through routing speed.
\end{caveat}

\end{document}
